\section{Encode a full record before adding temporal prediction}
A decoder receives bytes and constructs numeric display words. The encoder travels
in the other direction. It receives 4096 native \texttt{0x00RRGGBB} words and emits
one full-frame record: sixteen header bytes followed by 12\,288 RGB bytes. This
first stage deliberately emits a full frame at any valid sequence index. Full
frames are permitted between scheduled boundaries, so its records remain usable
when delta selection is added later.

Implement \texttt{video\_encode\_full}. Pass the four-byte-aligned source in
\texttt{\$a0}, its exact byte count 16\,384 in \texttt{\$a1}, the byte-addressed
record destination in \texttt{\$a2}, and a 64-byte, eight-byte-aligned native
context in \texttt{\$a3}. The destination need not be word aligned. The caller
supplies valid, disjoint source, destination, previous-image, context and private
stack windows. Retain the saved registers, frame pointer and stack pointer. Keep
the stack sixteen-byte aligned. The illustrated private frame is 96 bytes; an
equivalent implementation may use another layout.

\begin{table}[htbp]
\centering
\caption{Full-record encoder context. All fields occupy eight native bytes.}
\label{tab:video-full-encode-context}
\begin{tabular}{r p{1.6in} p{2.7in}}
\toprule
Offset & Field & Meaning \\
\midrule
0 & Destination capacity & Input, at least 12\,304 bytes. \\
8 & Next sequence & Input zero through 255; increment only on success. \\
16 & Full-frame period & Input one through 64; retain it. \\
24 & Previous-image address & Input pointer to four-byte-aligned storage. \\
32 & Previous-image capacity & Input, at least 16\,384 bytes. \\
40 & Last kind & Output zero for this full record. \\
48 & Encoded byte count & Output 12\,304. \\
56 & Retained suffix & Leave unchanged. \\
\bottomrule
\end{tabular}
\end{table}

Success returns zero in \texttt{\$v0} and 12\,304 in \texttt{\$v1}. It writes
the complete record, copies the source image to previous-image storage, increments
the sequence, and publishes the kind and encoded count. Retain the source on both
paths. Failure returns one and zero without storing to any destination,
previous-image or context byte. This includes invalid colors discovered at the
last pixel, not merely invalid capacities discovered at entry. Check counts,
sequence, period and native alignment before source or output access. Reject source
count/alignment and context alignment before reading context fields.

Every source word must have a zero high byte. Accepting a nonzero high byte and
discarding it would silently change the input contract. Inspect all 4096 words
before writing any record byte. Once validation and capacities succeed, the
remaining steps have admitted storage and fixed extents. This permits a direct
output pass without an extra payload buffer. The previous image changes only
after the complete record, including its checksum, has been written.

\begin{table}[htbp]
\centering
\caption{One numeric pixel has different memory bytes but the same RGB output.}
\begin{tabular}{l l l}
\toprule
Representation & Bytes for \texttt{0x00123456} & Interpretation \\
\midrule
Native little endian & \texttt{56 34 12 00} & Four-byte source word \\
Native big endian & \texttt{00 12 34 56} & The same numeric source word \\
MLV1 payload & \texttt{12 34 56} & Red, green, blue \\
\bottomrule
\end{tabular}
\end{table}

Load a native word numerically, extract channels with shifts, and store individual
bytes. Copying the first three source memory bytes would be wrong in both native
orders shown. The unused high byte occupies source storage but is not a color
sample. Its removal is a representation change, not temporal compression.

The default source is a deterministic gradient: pixel $i$ is
\texttt{((i * 257) XOR 0x123456) AND 0xFFFFFF}. Its first four payload triples are
\texttt{12 34 56}, \texttt{12 35 57}, \texttt{12 36 54} and \texttt{12 37 55};
the final triple is \texttt{02 3A A9}. At sequence zero, the completed header is
\texttt{00 00 00 00 00 00 00 00 00 00 30 00 F8 CD 76 19}.
The CRC covers header bytes zero through eleven followed immediately by the payload
for checksum purposes. Stored checksum bytes twelve through fifteen are excluded.

Attach this stage with
\texttt{python3 packages/034-lossless-video/construction/video-full-encode-construction/attach.py --source
seek-work.asm --output encoder-work.asm}. Earlier instructions remain intact under
the retained seek example entry. The new caller uses generated native pixels,
encodes a full record, then calls the existing record decoder on those bytes.
It needs no host video file. Preparation changes input words and context controls
without changing instructions. The fixture helper can also export the numeric
source words as JSON for independent inspection; it supplies no video encoder.
All output commands refuse an existing file.

Development status four can expose a partially written record after the header
regions are filled, or an updated previous image after its copy region is filled.
That is limited to unfinished construction. Completed failure one permits none
of those stores. Later delta and file-encoding stages will add prediction and
complete-file framing; a correct full record by itself is not yet an animation file.
